\section{Memory: del caché de contexto a una base de conocimiento sostenible}

\subsection{Diseño de memoria por capas}
El curso considera que Memory es el factor decisivo de la estabilidad del Agent. Sin memoria por capas, las tareas largas degeneran en «empezar de cero en cada turno». Esta clase recomienda distinguir al menos: working memory, episodic memory, semantic memory, procedural memory.

\begin{figure}[H]
\centering
\includegraphics[width=0.84\textwidth]{slide-04-memory-design.jpg}
\caption{Diapositiva del video en aproximadamente 00:52:40: el módulo de memoria no es una base vectorial única, sino un sistema jerárquico de lectura y escritura}
\end{figure}

\begin{table}[H]
\centering
\begin{tabular}{p{2.8cm}p{3.6cm}p{4.8cm}}
\toprule
\textbf{Capa de memoria} & \textbf{Contenido típico} & \textbf{Foco de ingeniería} \\
\midrule
Working memory & Estado del paso actual, variables temporales & Actualización rápida, baja latencia, recortable \\
Episodic memory & Trayectoria de la tarea, registros de fallas, ramas de decisión & Soporta la revisión posterior y la transferencia de estrategias \\
Semantic memory & Hechos estables, conocimiento del dominio, documentación de herramientas & Precisión de recuperación y consistencia de versiones \\
Procedural memory & Flujos de trabajo reutilizables, plantillas de estrategia, patrones de llamada a herramientas & Diseño de la granularidad de abstracción y las condiciones de activación \\
\bottomrule
\end{tabular}
\caption{División del trabajo de las cuatro capas de memoria en el sistema del Agent}
\end{table}

\begin{importantbox}{La esencia de Memory es la lectura y escritura controladas, no el almacenamiento en un solo sentido}
El sistema de memoria debe responder cuatro preguntas: qué escribir, cuándo escribir, cómo leer y cuándo olvidar. Si solo se hace «escritura total», pronto aparecerán acumulación de ruido y desajustes de recuperación, que terminan por deteriorar la calidad del razonamiento y la planificación.
\end{importantbox}

\begin{knowledgebox}{La API mínima de las operaciones de Memory}
\begin{itemize}
    \item `write(event, importance, scope)`: la escritura incluye la importancia y el alcance;
    \item `retrieve(query, budget, freshness)`: la recuperación fija de manera explícita el presupuesto y la restricción de vigencia;
    \item `compress(window)`: resume y comprime según la etapa de la tarea;
    \item `evict(policy)`: ejecuta la eliminación por caducidad y la fusión de conflictos.
\end{itemize}
\end{knowledgebox}

\begin{warningbox}{La contaminación de la memoria es la causa oculta más común de las fallas en tareas largas}
Si se recuperan de manera repetida conclusiones erróneas del historial, el Agent cometerá el mismo error de forma estable. La propagación de la contaminación debe reducirse mediante la confianza en la fuente, la detección de conflictos y mecanismos de reversión.
\end{warningbox}

\subsection{Resumen de la sección}
Memory determina si el Agent puede acumular capacidad efectiva a través de turnos y de tareas. Estar organizada por capas, ser controlable y ser reversible son los tres requisitos mínimos de un sistema de memoria de alta calidad.

